
\vspace{1cm}

\noindent  \textbf{Leia as instruções antes de fazer a lista:}

\vspace{0.3cm}
\noindent A resolução deve ser entregue em formato Jupyter Notebook ou Google Colab na plataforma Google Classroom, na tarefa designada e até o dia determinado.

\noindent  Não utilize laços de repetição, condicionais ou recursão nessa lista. Não utilize nenhum módulo, método ou função já existente no Python exceto pela função \code{print}. 

\vspace{0.5cm}


\begin{enumerate}[label=\textbf{Q\arabic*.}]
    \item Escreva uma função em Python que recebe por argumento o preço de um item do mercado em reais e o seu peso em kg e retorna a quantidade de reais por grama deste produto.

    \ifmostrarGabarito{
\begin{minted}{python}
# Resposta:
def calcular_preco_por_grama( preco_item_reais, peso_kg ):
    peso_g = peso_kg*1000
    return preco_item_reais / peso_g
\end{minted}
}


    \item Escreva uma função em Python que receba como argumento um número $x$ e que retorne um número $y$ dado pela função abaixo:
    \[
    y = x^{5/3} + \frac{x^2 - 8x}{10} - \left( \frac{5x^4 + 9x^3 - 10x^2}{x^5 + 2} \right)^{1/2}
    \]
    
    \ifmostrarGabarito
\begin{minted}{python}
# Resposta:
def calcular_y( x ):
    fator_a = x**(5/3)
    fator_b = (x**2 - 8*x)/10
    fator_c = ((5*(x**5) + 9*(x**3) - 10*(x**2))/(x**5 + 2))**(1/2)
    return fator_a + fator_b - fator_c
\end{minted}
\fi


    \item A física nos diz que é possível descrever o movimento de corpos quando estes se movem em linha reta e sem aceleração (isto é, com velocidade constante), de acordo com a seguinte fórmula:
    \[
    s = s_0 + v \cdot \Delta t
    \]
    onde $s$ é a posição final, $s_0$ é a posição inicial, $v$ é a velocidade em que o corpo está se movendo e $\Delta t$ é o intervalo de tempo decorrido para que objeto vá de $s_0$ a $s$ com essa velocidade. Este tipo de movimento é chamado de movimento retilíneo uniforme (MRU). Escreva uma função em Python que recebe por argumento a posição inicial ($s_0$), a posição final ($s$) e o intervalo de tempo ($\Delta t$) e retorna a velocidade de um corpo se movendo em MRU.

    \ifmostrarGabarito
\begin{minted}{python}
# Resposta:
def calcular_velocidade( s_0, s, delta_t ):
    return (s-s_0)/delta_t
\end{minted}
\fi



    \item Escreva uma função em Python que recebe uma temperatura em Kelvin e retorna o seu equivalente em Fahrenheit sabendo que:
    \[
    \frac{T_{\text{Fahrenheit}} - 32}{9} = \frac{T_{\text{Kelvin}} - 273}{5}
    \]

    \ifmostrarGabarito{
\begin{minted}{python}
# Resposta:
def calcular_fahrenheit( temp_kelvin ):
    return ((temp_kelvin - 274)*9)/5 + 32
\end{minted}
}


    \item Escreva uma função que calcule o perímetro de um retângulo sabendo que valores são passados por argumento: o valor da sua diagonal e de um dos lados (nessa ordem). Lembre-se que a diagonal do retângulo pode ser pensada como a hipotenusa do triângulo formado por dois lados adjacentes e a diagonal. Lembre-se também que não é permitido importar nenhum módulo ou usar funções prontas do Python. O cálculo de uma raiz quadrada pode ser feito elevando a 0.5.

    

\ifmostrarGabarito{
\begin{minted}{python}
# Resposta:
def calcular_perimetro_retangulo( diagonal, lado ):
    # a^2 + b^2 = c^2
    c = diagonal
    a = lado
    b = (c**2 - a**2)**(0.5)
    perimetro = 2*a + 2*b
    return perimetro
\end{minted}
}


    \item A média, a variância e o desvio padrão são medidas muito utilizadas em estatística. A média é um valor que demonstra a concentração dos dados de uma distribuição como um ponto de equilíbrio. Já a variância e o desvio padrão são usados para indicar o quanto um conjunto de dados se desvia da média. A diferença entre a variância e o desvio padrão é que a variância é expressa em unidades quadráticas, enquanto o desvio padrão permite uma interpretação da dispersão na mesma unidade dos dados originais. Considerando um conjunto com seis valores, a média, a variância e o desvio padrão desse conjunto são calculados da seguinte forma:
    \[
    m = \frac{x_1 + x_2 + x_3 + x_4 + x_5 + x_6}{6}
    \]
    \[
    V = \frac{(x_1 - m)^2 + (x_2 - m)^2 + (x_3 - m)^2 + (x_4 - m)^2 + (x_5 - m)^2 + (x_6 - m)^2}{6} = \sigma^2,
    \]
    onde $m$ é a média, $V$ é a variância e $\sigma$ é o desvio padrão. Escreva uma função em Python que recebe seis números por argumento, e calcula e retorna a média $m$, a variância $V$ e o desvio padrão $\sigma$, nesta ordem. Lembre-se que não é permitido importar nenhum módulo ou usar funções prontas do Python.

    \ifmostrarGabarito{
\begin{minted}{python}
# Resposta:
def calcular_media_e_desvio_padrao( x1, x2, x3, x4, x5, x6 ):
    media = (x1+x2+x3+x4+x5+x6)/6
    sigma2 = (x1-media)**2
    sigma2+= (x2-media)**2
    sigma2+= (x3-media)**2
    sigma2+= (x4-media)**2
    sigma2+= (x5-media)**2
    sigma2+= (x6-media)**2
    sigma2/=6
    sigma = sigma2**(1/2)
    return media, sigma2, sigma
\end{minted}
}

    
\end{enumerate}
